%% Đề 34 -- de-tuyen-sinh-lop-10-mon-toan-chuyen-nam-2024-2025-so-gddt-lang-son.pdf
\deheader{ĐỀ SỐ 34}%
  {Kỳ thi Tuyển sinh vào lớp 10 THPT — Năm học 2024--2025}%
  {Sở GD\&ĐT Lạng Sơn | Chỉ phần Đại số}
\addcontentsline{toc}{section}{Đề 34 — Tuyển sinh vào lớp 10 THPT — Năm học 2024--2025}

\begin{de_cau}{Câu 1}{2,0 điểm}
\begin{enumerate}[label=\alph*)]
  \item Tính: $A = \sqrt{3+2\sqrt{2}} - \sqrt{3-2\sqrt{2}}$.
  \item Rút gọn:
  \[B = \dfrac{x-1}{\sqrt{x}-1} - \dfrac{\sqrt{x}+1}{x+\sqrt{x}}\]
  với $x > 0,\; x \neq 1$.
\end{enumerate}
\end{de_cau}

\begin{loi_giai_cau}
\textbf{a)} $\sqrt{3+2\sqrt{2}}=\sqrt{(\sqrt{2}+1)^2}=\sqrt{2}+1$; $\sqrt{3-2\sqrt{2}}=\sqrt{2}-1$.

$A=(\sqrt{2}+1)-(\sqrt{2}-1)=\boxed{2}$.

\textbf{b)} $\dfrac{x-1}{\sqrt{x}-1}=\sqrt{x}+1$. $\dfrac{\sqrt{x}+1}{x+\sqrt{x}}=\dfrac{\sqrt{x}+1}{\sqrt{x}(\sqrt{x}+1)}=\dfrac{1}{\sqrt{x}}$.

$B=\sqrt{x}+1-\dfrac{1}{\sqrt{x}}=\boxed{\dfrac{x+\sqrt{x}-1}{\sqrt{x}}}$.
\end{loi_giai_cau}

\begin{de_cau}{Câu 2}{2,0 điểm}
\begin{enumerate}[label=\alph*)]
  \item Giải phương trình: $x^2 - 8x + 15 = \sqrt{x-3}$.
  \item Giải hệ phương trình:
  \[\begin{cases} x^2+y^2+2xy=9 \\ x^3+y^3=9 \end{cases}.\]
\end{enumerate}
\end{de_cau}

\begin{loi_giai_cau}
\textbf{a)} ĐKXĐ: $x\ge3$. Đặt $t=\sqrt{x-3}\ge0$, $x=t^2+3$:
\[(t^2+3)^2-8(t^2+3)+15=t \Rightarrow t^4-2t^2-t=0 \Rightarrow t(t^3-2t-1)=0 \Rightarrow t(t-1)(t^2+t+1)=0 \cdot(\text{không đúng}).\]
$t^3-2t-1=(t+1)(t^2-t-1)$: $t=0,t=-1$(loại),$t=\frac{1+\sqrt{5}}{2}$.

$t=0\Rightarrow x=3$ \checkmark ; $t=\frac{1+\sqrt{5}}{2}\Rightarrow x=3+\frac{3+\sqrt{5}}{2}$.

\textbf{b)} $(x+y)^2=9\Rightarrow x+y=\pm3$.

$x^3+y^3=(x+y)(x^2-xy+y^2)=(x+y)((x+y)^2-3xy)=9$.

$x+y=3$: $3(9-3xy)=9\Rightarrow 9-3xy=3\Rightarrow xy=2$. $x,y$ là nghiệm $t^2-3t+2=0\Rightarrow(t-1)(t-2)=0$: $(x;y)\in\{(1;2),(2;1)\}$.

$x+y=-3$: $-3(9-3xy)=9\Rightarrow 9-3xy=-3\Rightarrow xy=4$. $t^2+3t+4=0$: $\Delta<0$ -- vô nghiệm thực.

Nghiệm: $\boxed{(x;y)\in\{(1;2),(2;1)\}}$.
\end{loi_giai_cau}

\begin{de_cau}{Câu 3}{2,0 điểm}
\begin{enumerate}[label=\alph*)]
  \item Tìm các cặp số nguyên $(x;y)$ thỏa mãn: $x^2 + 2y^2 = 2x + 3$.
  \item Cho $p$ là số nguyên tố lớn hơn 3. Chứng minh $p^2-1$ chia hết cho $24$.
\end{enumerate}
\end{de_cau}

\begin{loi_giai_cau}
\textbf{a)} $(x-1)^2+2y^2=4$. $2y^2\le4\Rightarrow y\in\{-1,0,1\}$.

$y=0$: $(x-1)^2=4\Rightarrow x=3$ hoặc $x=-1$.
$y=\pm1$: $(x-1)^2=2$ -- không nguyên.

Nghiệm: $(x;y)\in\{(3;0),(-1;0)\}$.

\textbf{b)} $p>3$ nguyên tố nên $p$ lẻ và $p\not\vdots3$. $p^2-1=(p-1)(p+1)$.

$p$ lẻ $\Rightarrow p-1,p+1$ chẵn và là hai số chẵn liên tiếp $\Rightarrow(p-1)(p+1)\vdots8$.

$p\not\equiv0\pmod3\Rightarrow p\equiv\pm1\pmod3\Rightarrow p^2\equiv1\pmod3\Rightarrow p^2-1\vdots3$.

Vì $\gcd(8,3)=1$: $24\mid p^2-1$. $\square$
\end{loi_giai_cau}

\begin{de_cau}{Câu 5}{0,5 điểm}
Cho $a, b, c > 0$ thỏa mãn $abc = 1$. Chứng minh rằng:
\[\dfrac{a}{b+1} + \dfrac{b}{c+1} + \dfrac{c}{a+1} \ge \dfrac{3}{2}.\]
\end{de_cau}

\begin{loi_giai_cau}
Áp dụng AM-GM: $\dfrac{a}{b+1}\ge\dfrac{a}{b+1}$. Theo Nesbitt mở rộng và $abc=1$:

Dùng Cauchy-Schwarz: $\displaystyle\sum\dfrac{a}{b+1}\ge\dfrac{(a+b+c)^2}{\sum a(b+1)}=\dfrac{(a+b+c)^2}{ab+a+bc+b+ca+c}=\dfrac{(a+b+c)^2}{(ab+bc+ca)+(a+b+c)}$.

Đặt $s=a+b+c\ge3\sqrt[3]{abc}=3$, $q=ab+bc+ca\ge3\sqrt[3]{a^2b^2c^2}=3$.

$\displaystyle\sum\dfrac{a}{b+1}\ge\dfrac{s^2}{q+s}\ge\dfrac{9}{6}=\dfrac{3}{2}$ (khi $s=3,q=3$). $\square$
\end{loi_giai_cau}
